\section{Exact Dimension-Three Characterization}\label{sec:three}

We now specialize the simplex reduction to real \(3\times3\) matrices. Assume first that \(-A\) is a strict \(P\)-matrix and write
\[
p_i=-a_{ii}>0,\qquad
m_{12},m_{13},m_{23}>0,\qquad
q=-\det A>0,
\]
with the normalized invariants in Equation~\eqref{eq:four-invariants}. For
\[
D=\diag(d_1,d_2,d_3)\succ0
\]
the characteristic polynomial is
\begin{equation}
\det(\lambda I-DA)
=
\lambda^3+a_D\lambda^2+b_D\lambda+c_D,
\label{eq:cubicD}
\end{equation}
where
\begin{align}
a_D&=p_1d_1+p_2d_2+p_3d_3,\\
b_D&=m_{12}d_1d_2+m_{13}d_1d_3+m_{23}d_2d_3,\\
c_D&=q\,d_1d_2d_3.
\end{align}
With
\[
x_i=\frac{p_id_i}{a_D},
\qquad x\in\Delta_2^\circ,
\]
define
\begin{equation}
B_\beta(x)
=
\beta_{12}x_1x_2+
\beta_{13}x_1x_3+
\beta_{23}x_2x_3.
\label{eq:Bbeta}
\end{equation}
Then
\begin{equation}
\frac{b_D}{a_D^2}=B_\beta(x),
\qquad
\frac{c_D}{a_D^3}=\kappa x_1x_2x_3.
\label{eq:normalized-coeffs}
\end{equation}

\subsection{The Fixed-Cubic Matignon Boundary}

The fixed-polynomial component is a fractional Routh--Hurwitz/root-location problem; related coefficient-space criteria are established in the literature \cite{CermakNechvatal2017,BourafaAbdelouahabMoussaoui2020}. We state the exact form required for the diagonal-orbit elimination.

\begin{Lemma}[boundary value for a positive cubic]\label{lem:cubic-boundary}
Let
\[
p(\lambda)=\lambda^3+a\lambda^2+b\lambda+c,
\qquad a,b,c>0,
\]
and let
\[
\frac{\pi}{3}<\theta<\frac{\pi}{2}.
\]
Set
\[
u=\cos\theta,\qquad K=1-4u^2>0,
\]
and define
\begin{equation}
r_\theta(a,b)
=
\frac{ua+\sqrt{u^2a^2+Kb}}{K},
\qquad
H_\theta(a,b)
=
r_\theta(a,b)^2\bigl(a+2ur_\theta(a,b)\bigr).
\label{eq:Htheta}
\end{equation}
Then
\begin{equation}
|\arg\lambda_j|>\theta\quad\forall j
\quad\Longleftrightarrow\quad
c<H_\theta(a,b).
\label{eq:fixed-cubic-iff}
\end{equation}
Equality places one conjugate pair exactly on the rays \(|\arg\lambda|=\theta\). Moreover,
\[
H_\theta(sa,s^2b)=s^3H_\theta(a,b)
\qquad(s>0).
\]
\end{Lemma}

\begin{proof}
A polynomial with positive coefficients has no positive real zero, so loss of Matignon stability can occur only through a conjugate pair crossing a boundary ray. Put \(\lambda=re^{i\theta}\), \(r>0\), in \(p(\lambda)=0\). Dividing the imaginary part by \(r\sin\theta\) gives
\[
r^2\frac{\sin3\theta}{\sin\theta}
+
ar\frac{\sin2\theta}{\sin\theta}
+b=0.
\]
Using
\[
\frac{\sin3\theta}{\sin\theta}=4u^2-1=-K,
\qquad
\frac{\sin2\theta}{\sin\theta}=2u,
\]
we obtain
\[
Kr^2-2uar-b=0,
\]
whose unique positive solution is \(r=r_\theta(a,b)\). The real part then gives
\[
c=r^2(a+2ur)=H_\theta(a,b).
\]
For \(c>0\) sufficiently small, the roots consist of a small negative real root and a perturbation of the roots of \(\lambda^2+a\lambda+b\), hence are inside the Matignon region. For \(c\to\infty\), the roots approach the cube-root directions of \(-c\), including arguments \(\pm\pi/3\), which violate \(\theta>\pi/3\). Since there is only one boundary value, Equation~\eqref{eq:fixed-cubic-iff} follows. Homogeneity is immediate from Equation~\eqref{eq:Htheta}.
\end{proof}

For \(2/3<\alpha<1\), put \(\theta=\alpha\pi/2\) and define
\begin{equation}
h_\alpha(b):=H_\theta(1,b).
\end{equation}
Equivalently,
\begin{equation}
r_\alpha(b)
=
\frac{u+\sqrt{u^2+(1-4u^2)b}}{1-4u^2},
\qquad
h_\alpha(b)
=
r_\alpha(b)^2\bigl(1+2ur_\alpha(b)\bigr),
\label{eq:halpha}
\end{equation}
where \(u=\cos(\alpha\pi/2)\).

\subsection{Exact Elimination of the Positive-Diagonal Orbit}

Define
\begin{equation}
T_\alpha(\beta)
=
\min_{x\in\Delta_2^\circ}
\frac{h_\alpha(B_\beta(x))}{x_1x_2x_3}.
\label{eq:Talpha}
\end{equation}
The minimum is attained in the simplex interior because the numerator remains positive on the closed simplex while \(x_1x_2x_3\to0\) on its boundary.

\begin{Theorem}[exact three-dimensional Matignon \(D\)-stability]\label{thm:c10}
Let \(-A\) be a strict \(P\)-matrix and let \(2/3<\alpha<1\). Then
\begin{equation}
\boxed{
A\in\mathcal F_\alpha^{(3)}
\quad\Longleftrightarrow\quad
\kappa<T_\alpha(\beta).
}
\label{eq:c10-iff}
\end{equation}
\end{Theorem}

\begin{proof}
For a positive diagonal \(D\), let \(s=a_D\) and let \(x\in\Delta_2^\circ\) be its normalized orbit coordinate. Equations~\eqref{eq:normalized-coeffs} give
\[
a_D=s,\qquad
b_D=s^2B_\beta(x),\qquad
c_D=s^3\kappa x_1x_2x_3.
\]
By Lemma~\ref{lem:cubic-boundary} and homogeneity,
\[
\spec(DA)\subset\Sigma_\alpha
\]
is equivalent to
\[
\kappa
<
\frac{h_\alpha(B_\beta(x))}{x_1x_2x_3}.
\]
The normalization \(D\mapsto x\) is onto the open simplex modulo common scale. Hence the inequality must hold for every \(D\succ0\) if and only if it holds for every \(x\in\Delta_2^\circ\), which is exactly Equation~\eqref{eq:c10-iff}.
\end{proof}

The theorem is an explicit low-dimensional solution of a generalized multiplicative \(D\)-stability problem. Abstract forbidden-boundary criteria for conic regions and complements are already known \cite{KushelPavaniJDDE}; the step established in Theorem~\ref{thm:c10} is the complete elimination of the universal diagonal multiplier into the four invariants and the scalar threshold~\eqref{eq:Talpha}.

\subsection{Interior and the Classical Cain Boundary}

The strict-\(P\) hypothesis is exactly the full-dimensional stratum relevant to the interior.

\begin{Lemma}[principal-minor necessity]\label{lem:P0}
If \(A\in\mathcal F_\alpha^{(3)}\), then
\[
-a_{ii}\ge0,\qquad
\det A[\{i,j\}]\ge0,\qquad
-\det A>0.
\]
Consequently, every interior point of \(\mathcal F_\alpha^{(3)}\) has \(-A\) strict \(P\).
\end{Lemma}

\begin{proof}
If some \(a_{ii}>0\), let the other row scalings tend to zero; the limiting matrix has a positive real eigenvalue, and nearby positive diagonal scalings violate Matignon stability. If an order-two principal minor is negative, the same argument with the complementary row scaling tending to zero produces a \(2\times2\) block with a positive real eigenvalue. Finally a real \(3\times3\) Matignon-stable matrix has negative determinant, so \(-\det A>0\). Any zero order-one or order-two signed principal minor can be perturbed arbitrarily slightly to the forbidden sign, and therefore cannot occur at an interior point.
\end{proof}

At the classical endpoint \(\alpha=1\),
\[
h_1(b)=b,
\]
so
\begin{equation}
T_1(\beta)
=
\min_{x\in\Delta_2^\circ}
\left(
\frac{\beta_{23}}{x_1}
+
\frac{\beta_{13}}{x_2}
+
\frac{\beta_{12}}{x_3}
\right).
\end{equation}
Cauchy--Schwarz gives the exact value
\begin{equation}
\boxed{
T_1(\beta)
=
\left(
\sqrt{\beta_{12}}+
\sqrt{\beta_{13}}+
\sqrt{\beta_{23}}
\right)^2,
}
\label{eq:T1}
\end{equation}
which is Cain's strict-\(P\) real \(3\times3\) \(D\)-stability threshold \cite{Cain1976}.

For \(2/3<\alpha<1\), \(T_\alpha\) lies strictly above this classical surface. Indeed, from
\[
b=(1-4u^2)r^2-2ur
\]
and Equation~\eqref{eq:halpha},
\[
h_\alpha(b)-b
=
4u^2r^2+2ur(r^2+1)>0.
\]
Therefore
\begin{equation}
\boxed{
T_\alpha(\beta)>T_1(\beta)
\qquad(2/3<\alpha<1).
}
\label{eq:strict-gap}
\end{equation}

\begin{Corollary}[exact genuinely fractional band]\label{cor:fracband}
For \(2/3<\alpha<1\),
\begin{equation}
\operatorname{int}\mathcal P_\alpha^{(3)}
=
\left\{
A:
-A\text{ strict }P,\quad
T_1(\beta)<\kappa<T_\alpha(\beta)
\right\}.
\label{eq:fractional-band}
\end{equation}
For \(0<\alpha\le2/3\),
\begin{equation}
\operatorname{int}\mathcal F_\alpha^{(3)}
=
\{A:-A\text{ strict }P\},
\label{eq:low-order-F}
\end{equation}
and
\begin{equation}
\operatorname{int}\mathcal P_\alpha^{(3)}
=
\left\{
A:
-A\text{ strict }P,\quad
\kappa>T_1(\beta)
\right\}.
\label{eq:low-order-P}
\end{equation}
\end{Corollary}

\begin{proof}
The high-order statement follows from Theorem~\ref{thm:c10}, Equation~\eqref{eq:T1}, and strictness of the inequalities. In the low-order regime, Kellogg's eigenvalue wedge theorem for \(P\)-matrices \cite{Kellogg1972} implies that if \(-A\) is strict \(P\), every eigenvalue of \(-DA\) lies in
\[
|\arg\mu|<\pi-\frac{\pi}{3}=\frac{2\pi}{3}.
\]
After multiplication by \(-1\), the eigenvalues of \(DA\) satisfy the Matignon inequality for every \(\alpha\le2/3\). Lemma~\ref{lem:P0} gives necessity for interior points, while Cain's threshold separates classical \(D\)-stability from its complement.
\end{proof}

Figure~\ref{fig:c10geometry} shows the exact threshold geometry on symmetric and asymmetric slices. The blue Matignon boundary remains strictly above the orange Cain boundary, and the region between them is precisely the all-positive-diagonal fractional-only band.

\begin{figure}[!htbp]
\centering
\includegraphics[width=0.92\textwidth]{fig03_threshold_geometry.pdf}
\caption{Exact C-10 threshold geometry. The orange curve is Cain's classical boundary \(T_1(\beta)\); the blue curve is the exact Matignon boundary \(T_\alpha(\beta)\). Their strict separation creates the genuinely fractional band. Certified anchor markers are computational checks of exact formulas, not substitutes for Theorem~\ref{thm:c10}.}
\label{fig:c10geometry}
\end{figure}

\subsection{Minimum Dimension for a Robust Fractional-Only Class}

The two-dimensional obstruction and the three-dimensional band give the dimension result.

\begin{Theorem}[minimum robust dimension]\label{thm:mindim}
For every \(0<\alpha<1\),
\begin{equation}
\boxed{
\min\left\{
n:
\operatorname{int}
\left(
\mathcal F_\alpha^{(n)}
\setminus
\mathcal D_H^{(n)}
\right)\ne\varnothing
\right\}=3.
}
\label{eq:min-dim}
\end{equation}
\end{Theorem}

\begin{proof}
Dimension one is trivial and Theorem~\ref{thm:2d} gives empty full-dimensional interior in dimension two. For \(2/3<\alpha<1\), Equation~\eqref{eq:strict-gap} makes the interval
\[
T_1(\beta)<\kappa<T_\alpha(\beta)
\]
nonempty for every positive \(\beta\); any strict interior point is stable under small full-matrix perturbations. For \(0<\alpha\le2/3\), choose any strict-\(P(-A)\) matrix violating Cain's inequality \(\kappa<T_1(\beta)\). Corollary~\ref{cor:fracband} then places it in \(\mathcal P_\alpha^{(3)}\), and strictness again gives a full-dimensional neighborhood.
\end{proof}
